\documentclass[11pt]{article}
\usepackage[a4paper,margin=26mm]{geometry}
\usepackage{amsmath,amssymb,amsthm,hyperref}
\setlength{\emergencystretch}{3em}
\newtheorem{theorem}{Theorem}
\newtheorem{proposition}[theorem]{Proposition}
\newcommand{\OO}{\mathcal O}
\newcommand{\MA}{\operatorname{MA}}
\newcommand{\vol}{\operatorname{vol}}
\title{Horizontal fixed factors and the vertical term in moving models}
\author{Unreviewed correction, continuation 209}
\date{6 October 2026}
\begin{document}
\maketitle
\begin{abstract}
Removing the generic fixed divisor from a time model requires retaining
its pulled vertical part. We give the exact full-system formula and an
explicit big nonnef surface example. Omitting that term changes a zero
auxiliary metric into a nontrivial metric, changes its energy from zero
to $-1/3$, and moves its unit mass from the generic point to a divisorial
point. This corrects a geometric step in earlier unreviewed drafts; it
does not certify their entire equation theorem.
\end{abstract}

\section*{The actual full-system formula}
We use only finite actual model intersections. For a nef compactified
model $C$ with generic bundle $A$ and twist $cF$, its raw energy is
$C^{n+1}/(n+1)-cA^n$. Its raw measure has coefficient
$b_Q(C|_Q)^n$ at each central prime $Q$ of multiplicity $b_Q$.
These are the finite intersection conventions; no unbounded operator
or nonample equation theorem is used below.
Let $X$ be smooth projective over an algebraically closed field of
characteristic zero, $L$ an actual big rational bundle, and $m$ a cleared
degree with nonempty complete system. Resolve its base ideal once:
\[
 p_m^*L=A_m+T_m,\qquad mA_m\text{ generated},\quad mT_m\geq0.
\]
Fix a smooth equivariant model carrying a rational vertical divisor $D$
over $X\times\mathbb P^1$. On a common smooth determination mapping
to the actual generic product $X_m\times\mathbb P^1$, write
\[
 \rho^*T_m=T_m^{\rm hor}+T_m^{\rm ver}.
                                                        \tag{1}
\]
Here the first term is the sum of horizontal prime components and the
second is supported over zero. Both are actual divisors; they are not
interchangeable numerical representatives. Let $F$ be the full fiber.
For a sufficiently large rational $c$ and a cleared degree $m$, the
complete time lattice of $\alpha_D(c)=\rho^*p_m^*L+D+cF$ has the
entire absolute basis
\[
 t^{-(w_j+mc)+r}s_j,\qquad 0\leq r\leq w_j+mc,
                                                        \tag{2}
\]
where $(s_j)$ is an adapted basis of the whole original $H^0(X,mL)$.
We may refine to resolve every base ideal below. All formulas are pulled
to that common determination.

\begin{theorem}[Retained vertical fixed part]\label{formula}
Divide the sections in (2) by the canonical section of
$mT_m^{\rm hor}$, and let $G_{m,D}^{\rm ver}\geq0$ be the
resolved residual base divisor. The actual complete moving bundle is
\[
 C_{m,D}=\rho^*A_m+D+T_m^{\rm ver}+cF
                       -m^{-1}G_{m,D}^{\rm ver}.          \tag{3}
\]
It is globally generated in degree $m$, and hence nef. Its generic
bundle is $A_m$. The bounded auxiliary quotient scalar is the model
scalar represented by $D+T_m^{\rm ver}-m^{-1}G_{m,D}^{\rm ver}$.
With $q_{T_m}$ the actual canonical divisor scalar, its physical
original stage satisfies $\phi_{m,D}=q_{T_m}+\psi_{m,D}$.
Formula (3) is the moving part of the \emph{full} absolute system (2).
\end{theorem}
\begin{proof}
Multiplication by the fixed section identifies
$H^0(X_m,mA_m)$ with $H^0(X,mL)$: every original section has exactly
the prescribed generic base-divisor factor, and normal proper
birational descent identifies the remaining spaces. Every section of
(2) therefore vanishes along $mT_m^{\rm hor}$. This gives an actual
division of regular sections and puts the residual sections in
\[
 m\alpha_D(c)-mT_m^{\rm hor}
   =m(\rho^*A_m+D+T_m^{\rm ver}+cF).                      \tag{4}
\]
Away from zero they generate that generic bundle: all time monomials
are present and the original generic moving system is complete and
generated. At infinity the endpoint monomials also recover the entire
generic generated system. Consequently the residual base ideal is
supported at zero. Principalization makes it the invertible ideal of
the effective vertical $G_{m,D}^{\rm ver}$. Removing that ideal gives
(3), generated by the same full list of sections. It removes every
fixed component of the entire absolute system, exactly once.

The generic residual sections are $s_j/s_{mT_m}$. Taking their
weighted maximum gives the actual bounded generated scalar on $A_m$.
Multiplying back gives the physical identity at every point where its
divisor factor is finite, and gives the same literal $-\infty$ value
at supported zeros. The resulting model scalar has central divisor
exactly the one in (3) minus its fiber twist. Thus the scalar and actual
bundle identifications agree, without subtracting infinite values.

One cannot replace the division in (4) by division by the full pulled
$mT_m$. The obstacle divisor $D$ may already consume some vertical
vanishing of the generic factor. Those components then cease to be
fixed components of the absolute section system. The example below
makes this failure explicit.
\end{proof}

For each fixed obstacle on the original total model, (3) remains its
complete moving bundle. Hence the usual complete-system lower positive
product convergence applies to this $C_{m,D}$ and the fixed
$\alpha_D(c)$, rather than to the formula with $T_m^{\rm ver}$ omitted.
On a model where $T_m^{\rm ver}=0$ the shorter formula is valid.
The correction is necessary after arbitrary common refinements, which
are needed to compare several obstacles and fixed model directions.

\section*{A complete big nonnef surface computation}
Let $X=\operatorname{Bl}_p\mathbb P^2$, let $H$ be the pullback of
$\OO_{\mathbb P^2}(1)$ and let $E$ be the exceptional curve. Put
\[
 L=H+E,\qquad H^2=1,\quad H E=0,\quad E^2=-1.
\]
For every $m\geq1$, restriction to $E$ shows that each section of
$mL$ vanishes along $E$: its restricted degree is negative. Repeating
after each removal gives
\[
 H^0(X,mL)=s_E^m H^0(X,mH).
                                                        \tag{5}
\]
Thus the whole moving pair is $A_m=H$, $T_m=E$, and $\vol(L)=1$.
Write $q_E(v)=-v(E)$ in the canonical norm convention.

Blow up $Z=E\times\{0\}$ in $X\times\mathbb P^1$, and call the
result $\mathcal Y$, its exceptional divisor $B$, and the horizontal
strict transform of $E\times\mathbb P^1$ by $E^{\rm hor}$. Then
\[
 \rho^*E=E^{\rm hor}+B.                                \tag{6}
\]
Take the original obstacle divisor $D=-B$. The blowup ideal is
$(\OO(-E),t)$; its normalized model scalar is
\[
 f_D(v)=-\min(v(E),1),\qquad q_E\leq f_D\leq0.         \tag{7}
\]
In (5), every section satisfies $\log|s|\leq mq_E\leq mf_D$,
whereas evaluation at the generic trivial valuation bounds every
feasible weight by zero. Thus the complete original flag has every
weight equal to zero. The physical complete stage is exactly $q_E$,
and its auxiliary quotient on $H$ is exactly zero in every degree.
In particular its auxiliary energy is zero and its raw measure is
$\delta_{v_0}$.

The original complete section trace is also exactly zero. Thus the
contradiction below concerns the actual finite flag and its physical
quotient before any envelope theorem or unbounded energy extension
is invoked.

The absolute class is literally
\[
 \alpha_D(c)=\rho^*L-B+cF
            =\rho^*H+E^{\rm hor}+cF.                   \tag{8}
\]
For any integer $c\geq1$ its complete sections after division by
$s_{E^{\rm hor}}^m$ are all the sections of
$m(\rho^*H+cF)$, pulled from the product. This bundle is generated;
the residual vertical base divisor is zero. Formula (3) therefore gives
\[
 C_{m,D}=\rho^*H-B+B+cF=\rho^*H+cF.                    \tag{9}
\]
This is exactly the zero auxiliary quotient just computed. Removing
the full pulled $E$ would wrongly retain the negative $-B$.

\begin{proposition}[The omitted term changes energy and measure]
The class $C'=\rho^*H-B+cF$ is a different generated nef model for
$c\geq1$. Its auxiliary scalar is $-\min(v(E),1)$ on $H$, its raw
measure is $\delta_{\operatorname{ord}_E}$ and its raw energy is
$-1/3$. Thus it cannot represent the zero quotient in (9).
\end{proposition}
\begin{proof}
The two-dimensional pencil of plane lines through $p$ generates
$H-E$, and all three plane line sections generate $H$. Give the
pencil weight zero and the remaining generator weight $-1$.
The associated relative base ideal is $(\OO(-E),t)$; blowup resolves
it and its model polarization is precisely $\rho^*H-B+cF$.
The absolute time systems are generated for $c\geq1$: the two
time endpoints recover the pencil and all three sections, and the
full time monomials generate between them. Equivalently the ideal
is generated by the pencil and the $t$-weighted complete system.

Let $X'_0$ be the central strict transform. Blowing up the Cartier
$E$ in $X$ leaves $X'_0\simeq X$. The full fiber is $X'_0+B$, both
components having multiplicity one. Restricting $C'$ to $X'_0$
gives $H-E$, of square zero. On $B$, the pullback of $H$ is numerically
trivial, as $H|_E$ has degree zero. The normal bundle of the blowup
center is $N=\OO_E(-1)\oplus\OO_E$. With the quotient convention
$B=\mathbb P(N^\vee)$, set $\zeta=c_1(\OO_B(1))$; then
$\OO_B(B)=\OO_B(-1)$ and $\int_B\zeta^2=\deg N^\vee=1$.
Hence $C'|_B=\zeta$ numerically and has square one.
The central valuations are respectively $v_0$ and $\operatorname{ord}_E$.
The exact component intersection formula yields
\[
 \MA_{H,\mathrm{raw}}(-\min(v(E),1))
          =0\delta_{v_0}+1\delta_{\operatorname{ord}_E}. \tag{10}
\]

The ordinary blowup identities also compute the energy directly:
\[
 (\rho^*H+cF)^3=3c,\quad
 (\rho^*H+cF)^2B=0,\quad
 (\rho^*H+cF)B^2=-H E=0,\quad B^3=1.
\]
Therefore $C'^3=3c-1$ and
\[
 e_H(-\min(v(E),1))=\frac{C'^3}{3}-cH^2=-\frac13.     \tag{11}
\]
The zero auxiliary scalar instead has $C^3=3c$, raw energy zero
and raw measure $\delta_{v_0}$. Every value is finite at the two
atoms; no supported-pole subtraction is involved.
\end{proof}

\section*{Application and limits of this correction}
The formula missing $T_m^{\rm ver}$ occurs in the unreviewed
round-163 manuscript (5.2), and the instruction to remove the entire
pulled factor occurs in the unreviewed continuation-208 manuscript (7).
Both frozen originals are preserved. This calculation establishes a
specific obstruction to taking those formulas literally on every
required common model. The restricted-lattice formula retaining the
vertical term, for example (9.2) of round 163, is compatible with (3).
If a proof uses the shorter formula only on determinations where its
vertical term vanishes, it must explicitly show that this restriction
survives every required comparison and refinement.

The exact correction (3) preserves the complete moving-system
interpretation, so it provides a repair candidate for the energy and
measure limits. It does not by itself prove all subsequent unbounded
limits, mixed identities, comparison or attainment. Those require a
fresh whole-proof check with (3) substituted throughout. No active
review receives this new result as feedback. This note has no
independent acceptance and closes no original problem.
\end{document}
